\templateentry
{Base}
{}
{Base inteira do template de geometria com funções essenciais que vão ser usadas pela maioria dos outros templates. Esta versão usa \texttt{ll}, mas para alguns problemas, pode ser necessário trocar para \texttt{ld}. Na função \texttt{sgn}, se for usar com \texttt{ll}, troque \texttt{EPS} por 0.}
{geometry/core.cpp}

\templateentry
{Interseção de retas}
{\texttt{lineInter(a, b, c, d)} para retas dadas por dois pontos; $O(1)$.}
{Este snippet funciona com \texttt{ld}. Dados dois pontos de cada reta, \texttt{a,b} e \texttt{c,d}, ele encontra o ponto de interseção das retas infinitas \texttt{ab} e \texttt{cd}. Não é teste de segmentos. O snippet assume que as retas não são paralelas nem coincidentes, trate esse caso antes da chamada. Complexidade: $O(1)$.}
{geometry/line-intersection.cpp}

\templateentry
{Interseção de segmentos}
{\texttt{segInter} e \texttt{properInter} para dois segmentos, em $O(1)$.}
{Determina se dois segmentos se intersectam. Lembre-se de adaptar o \texttt{sgn} dependendo se vai usar \texttt{ll} ou \texttt{ld}. \texttt{properInter} testa apenas interseção estrita, excluindo os extremos dos segmentos. Complexidade: $O(1)$.}
{geometry/segment-intersection.cpp}

\templateentry
{Área de um polígono}
{\texttt{polyArea(p)} retorna a área orientada multiplicada por 2, em $O(n)$.}
{Encontra a área orientada de um polígono, positiva se o sentido é anti-horário, e negativa se é horário. Esta versão encontra a área inteira no caso de inteiros (2x a área real), então é necessário dividir por 2 se quiser encontrar a área real do polígono. Complexidade: $O(n)$.}
{geometry/polygon-area.cpp}

\templateentry
{Verificar se um ponto está dentro de um polígono convexo}
{Point-in-convex-polygon em $O(\log n)$.}
{\texttt{pointInPoly(pt, q, border)} determina se o ponto \texttt{q} está dentro do polígono convexo \texttt{pt} no sentido anti-horário. Use \texttt{border = true} se quiser incluir pontos na borda do polígono também. Complexidade: $O(\log n)$.}
{geometry/point-in-polygon.cpp}

\templateentry
{Convex hull}
{\texttt{convHull(pt)} retorna os índices dos vértices do hull em $O(n \log n)$.}
{Encontra o menor polígono convexo que contém todos os pontos. Dado o vetor de pontos \texttt{pt}, encontra a lista de índices dos vértices do convex hull em sentido anti-horário, referenciando o vetor original. Por padrão, pontos colineares sobre as arestas não entram no hull; o comentário no snippet mostra o que trocar se quiser incluí-los. Complexidade: $O(n \log n)$.}
{geometry/convex-hull.cpp}

\templateentry
{Menor distância euclideana}
{Closest pair: retorna um par de índices em $O(n \log n)$.}
{Dada uma lista de pontos distintos, encontra dois pontos com a menor distância euclideana. Considere o caso de pontos repetidos fora da função. Complexidade: $O(n \log n)$.}
{geometry/closest-pair.cpp}

\templateentry
{Interseção de círculo e reta}
{\texttt{circLineInter(c, r, a, d)} retorna 0, 1 ou 2 pontos, em $O(1)$.}
{Encontra a interseção de um círculo de centro \texttt{c} e raio \texttt{r}, e uma reta que passa por \texttt{a} e tem direção \texttt{d}. O retorno é a lista dos pontos de interseção entre o círculo e a reta infinita; ela vem com tamanho $0$, $1$ ou $2$. Complexidade: $O(1)$.}
{geometry/circle-line.cpp}

\templateentry
{Interseção de círculos}
{\texttt{circCircInter(c1, r1, c2, r2)} retorna 0, 1 ou 2 pontos, em $O(1)$.}
{Encontra os pontos de interseção de dois círculos. O caso degenerado de círculos coincidentes deve ser tratado fora da função. Complexidade: $O(1)$.}
{geometry/circle-circle.cpp}

\templateentry
{Tangentes comuns a dois círculos}
{\texttt{circTangents(c1, r1, c2, r2)} retorna pares de tangência, em $O(1)$.}
{Encontra as tangentes comuns a dois círculos. retorna uma lista de pares \texttt{\{p1, p2\}}, com um ponto de tangência em cada círculo. O snippet encontra tangentes externas e internas, com até $4$ respostas. Casos degenerados como círculos coincidentes devem ser tratados fora; em casos tangentes pode ser necessário remover duplicatas se isso importar no problema. Complexidade: $O(1)$.}
{geometry/circle-tangents.cpp}

\templateentry
{Menor cobertura circular}
{Encontra o menor círculo que contém todos os pontos.}
{Encontra o menor círculo que contém todos os pontos.}
{geometry/enclosing-circle.cpp}

\templateentry
{Soma Minkowski de polígonos convexos}
{\texttt{minkowski(a, b)} para convexos em sentido anti-horário, em $O(|a|+|b|)$.}
{A soma minkowski de dois conjuntos de pontos é o conjunto das somas de todos os pares de pontos dos conjuntos: $A \oplus B = \{P = a + b; a \in A, b \in B\}$. Passe dois polígonos convexos em sentido anti-horário, sem repetir o primeiro vértice no fim; o retorno é o polígono $a \oplus b$, também em sentido anti-horário. A rotina pode manter vértices redundantes se houver arestas colineares; normalize depois se precisar. Complexidade: $O(|a| + |b|)$.}
{geometry/minkowski.cpp}

\templateentry
{Corte de Polígono}
{Corte de polígono por semiplano em $O(n)$.}
{\texttt{polyCut(p, a, b)} corta o polígono \texttt{p} pelo semiplano à esquerda da reta direcionada $a \to b$. A reta é infinita e o resultado sai em ordem circular. Complexidade: $O(n)$.}
{geometry/polygon-cut.cpp}
